\section{Introduction}

The Allen--Cahn equation~\cite{AllenCahn}
\begin{equation}\label{eq:AC-intro}
u_t = \varepsilon^2 \Delta u - f(u), \qquad f(u) = u^3-u = F'(u), \qquad F(u)=\tfrac14(u^2-1)^2,
\end{equation}
is a prototypical reaction--diffusion gradient flow describing phase separation in binary alloys, and the archetypal example used to benchmark structure-preserving numerical schemes for Ginzburg--Landau phase-field models~\cite{Du_review}. Two structural properties make it a natural test bed for a journal devoted to symmetry. First, the double-well potential $F$ is invariant under $u\mapsto-u$, which maps solutions to solutions; independently of this, the variational structure $u_t=-\delta E/\delta u$ of the free energy $E[u]=\int_\Omega(\frac{\varepsilon^2}{2}|\nabla u|^2+F(u))\,dx$ gives the energy-dissipation law $dE/dt=-\|u_t\|^2_{L^2}\le0$, and the sign condition $f(-1)\le0\le f(1)$ gives, by the comparison principle, the maximum bound principle $|u_0|\le1\Rightarrow|u(\cdot,t)|\le1$~\cite{ref_Du2021b}. Second, on a domain with reflection or rotational symmetry, a symmetric initial condition produces a solution that stays symmetric, because the isotropic Laplacian and the odd nonlinearity commute with the symmetry group. Numerical methods that do not respect these properties can produce spurious energy growth, solutions outside $[-1,1]$ or broken symmetry, which is why stability in energy norms is central to the design of discrete schemes~\cite{Baishemirov_Symmetry,Baigereyev_Math,Shiyapov2026}. Classical schemes with unconditional energy stability (convex splitting~\cite{Eyre}, scalar auxiliary variable (SAV) approaches~\cite{Shen_SAV,Huang_GSAV,Jiang_RSAV}) and with a discrete maximum principle~\cite{ref_Tang2016,Zhang_ETD2024,Gao_pIFRK2023,Du_MPPDG2022,Feng_MPPFD2021,Kim_MPP3D2025} are well developed. Energy decay alone does not imply $|u|\le1$, so admissibility has to be checked separately; we do so for every network trained on the one-dimensional forward problem.

Physics-informed neural networks (PINNs)~\cite{Raissi} approximate the solution of a \rev{partial differential equation (PDE)} by a network trained to minimize the residual at collocation points together with initial- and boundary-condition losses. They are mesh-free and compatible with data assimilation and inverse problems~\cite{Karniadakis2021,Yang_PINN,Lu_DeepXDE,Pateras2023}. For the Allen--Cahn equation, however, vanilla PINNs have documented difficulties. For small $\varepsilon$ the solution develops an interface of width $O(\varepsilon)$ that the spectral bias of fully connected networks makes hard to resolve~\cite{Wight,ref_Rahaman2019,Wang_NTK2022}; adaptive sampling~\cite{Wight}, causal~\cite{Wang_causal} and sequential~\cite{Mattey} training and Fourier-feature encodings~\cite{Liu_FourierBias2025} have been proposed as remedies. More generally, PINN training is hampered by loss-term imbalance and ill-conditioned differential operators~\cite{ref_Wang2021a,ref_Krishnapriyan2021,ref_Rathore2024}. Finally, the standard loss enforces only the pointwise residual and gives no structural guarantee that the energy decreases or that $|u|\le1$ holds; Sections~\ref{sec:results-energy} and~\ref{sec:results-balance} examine whether this matters in practice.

\rev{The closest relatives of this study place a discretization inside the network or add structure to the loss. Networks whose only input is time and whose output is the vector of coefficients of a spectral basis, with spatial derivatives taken exactly from the basis, were introduced by Xia, B\"ottcher and Chou~\cite{ref_Xia2023} (Chebyshev, Laguerre and Hermite bases for linear transport, heat and Schr\"odinger equations, including the estimation of a diffusivity). Spectral informed neural networks (SINNs)~\cite{ref_Yu2026SINN} take the frequency and time as input and output the corresponding coefficient of a Fourier basis, so that spatial derivatives become multiplications. Further coefficient-space networks are neural spectral methods for parametric PDEs~\cite{ref_Du2024} and the neuro-spectral architectures of~\cite{ref_Bizzi2025}, in which boundary conditions hold by the choice of basis; neural Galerkin schemes~\cite{ref_Bruna2024} and evolutional networks~\cite{ref_Du2021a} march network parameters in time instead. For the Allen--Cahn equation, K\"ut\"uk and Y\"ucel~\cite{ref_Kutuk2025} add a one-sided penalty on $\partial E/\partial t$ to a coordinate-based PINN, an energy-based sequential training is proposed in~\cite{Guo_EnergySeq2024}, dissipation of a modified energy holds by construction in networks built on the SAV approach~\cite{ref_Zhang2024,ref_Lin2025}, and the maximum bound is enforced by a limiter in a learned solver~\cite{ref_Geng2025}. Imposing boundary conditions, periodicity or conservation laws exactly through the architecture is well established~\cite{ref_Dong2021,ref_Sukumar2022,ref_Lu2021,ref_CardosoBihlo2025}, and for Chebyshev polynomials the Neumann-compatible basis is the classical one of the Galerkin method~\cite{Shen1995}. For forward problems, classical solvers are faster and at least as accurate as PINNs, including for the one-dimensional Allen--Cahn equation~\cite{ref_Grossmann2024,ref_McGreivy2024}.}

Symmetry has also been built into PINNs through the loss, the data or the architecture: invariant-surface conditions from Lie point symmetries~\cite{Zhang_SPINN2023}, symmetry-generated training data~\cite{Zhang_SymGroup2023,Brandstetter2022}, invariants and moving frames~\cite{Arora2024}, symmetry-based domain decomposition~\cite{Liu_SymDD2024}, and, closest in spirit to the present work, symmetry-preserving PINNs for nonlinear lattices, which enforce spatio-temporal parity by a group-equivariant architecture~\cite{Zhu_SPINN2022}. The Fourier-basis coefficient network of Section~\ref{sec:2d} (Fourier-Spectral-PINN) instead enforces the reflection symmetries $x\mapsto-x$ and $y\mapsto-y$ through the \emph{basis functions} of the output representation, so that no loss term is responsible for them. This is a special case of the hard-constraint design cited above, and we claim only its demonstration for the two-dimensional Allen--Cahn equation.

\textbf{Scope and contributions.} \rev{We do not claim that coefficient-space networks or energy penalties are new. This article is a careful, reproducible evaluation of a Chebyshev-coefficient PINN for the Allen--Cahn equation (Spectral-PINN) against strong baselines. Its contributions are: (i)~the use of the classical Neumann-compatible Chebyshev--Galerkin basis~\cite{Shen1995} as the output basis of the coefficient network, and the observation that a soft Neumann penalty fails at large mode number $K$ in both of its natural scalings (it vanishes when normalized with the mode number, which produced a spurious loss of accuracy in the first version of this study, and makes the loss ill-conditioned when it is not; Section~\ref{sec:results-k}); (ii)~a forward-problem study on three benchmarks (B1, reaction-dominated; B2, diffusion-active; B3, with sharp interfaces, on which the PINNs fail), with up to five seeds, a certified reference solution and the classical Chebyshev collocation solver as the strong baseline, and maximum-bound and energy diagnostics for every network (Sections~\ref{sec:results-forward}, \ref{sec:results-energy} and~\ref{sec:results-balance}); (iii)~an a-posteriori energy balance, proved in Section~\ref{sec:balance}, valid for \emph{any} trained trajectory, which bounds the energy growth of the learned solution by its PDE residual and boundary defect, both measured for every Spectral-PINN (Section~\ref{sec:results-balance}); (iv)~a controlled ablation of the energy penalty, which leaves the accuracy essentially unchanged in all regimes tested except training without warm-up, where it degrades it, and shows that monotone energy is not a diagnostic of accuracy here (Section~\ref{sec:results-energy}); (v)~a diagnosis of why joint training of the network and of $\varepsilon$ returns an estimate set by the starting value, an estimator without this defect (accurate on noise-free data, biased under noise), and a comparison with classical least squares including identifiability (Section~\ref{sec:results-inverse}); (vi)~the demonstration that a restricted Fourier basis preserves the reflection symmetries of the two-dimensional solution to round-off (Section~\ref{sec:2d-results}). The method is not offered as a replacement for classical solvers; Section~\ref{sec:niche} states what we think its niche is.}

This article extends a conference abstract~\cite{ICCMA2026abstract} that outlined a spectral PINN for the Allen--Cahn equation with an SAV formulation and reported no numerical results; the SAV formulation is not used here (we study the energy penalty of Section~\ref{sec:energy-loss} instead), and the formulation, the two-dimensional extension, the multi-seed experiments and the energy balance of Section~\ref{sec:balance} are new.
